\documentclass[12pt]{article}
\usepackage{amsmath, amssymb, amsthm}
\usepackage[margin=1in]{geometry}
\usepackage{hyperref}

\title{Gate 5: The Complete Causal Chain}
\author{Ryan O. Van Gelder \\ Citizen Gardens / Multiplicity Theory Framework}
\date{March 2026}

\begin{document}

\maketitle

\begin{abstract}
Gate 5 closes the five-gate validation sequence of the CEQG-RG-Langevin programme by establishing an unbroken, unidirectional causal chain from microscopic quantum geometry (AL-GFT or Track-B/C variants) through renormalization-group flow, stochastic influence functional, and Einstein–Langevin dynamics to the two target late-time observables: the effective local trispectrum amplitude \(g_{\rm NL}^{\rm eff}\) (CMB scales) and the structure-growth parameter \(S_8\) (LSS scales). The chain produces the non-tunable exponential correlation of Gate 3
\[
\log\left(\frac{g_{\rm NL}}{g_{\rm NL0}}\right) = \frac{1}{c \cdot n_{\rm NG}} \cdot \frac{\Delta S_8}{S_8} + \mathcal{O}\left(\left(\frac{\Delta S_8}{S_8}\right)^2\right),
\]
without extraneous parameters. With all acceptance criteria from Gates 1–4 satisfied, the framework is fully specified, truncation-controlled, and ready for joint Bayesian forecasting and confrontation with forthcoming DESI Y5 + LiteBIRD + Euclid data.
\end{abstract}

\section{Requirement Statement}
Gate 5 requires demonstration of a complete, modularly validated causal pathway free of ad-hoc insertions:
\[
\begin{array}{c}
\text{Microscopic UV boundary conditions (Gate 1)} \\
\downarrow \\
\text{RG-derived cosmological prior (Gate 2)} \\
\downarrow \\
\text{Controlled truncation (Gate 4)} \\
\downarrow \\
\text{Stochastic noise kernel and running vacuum} \\
\downarrow \\
\text{Modified Einstein–Langevin dynamics} \\
\downarrow \\
\text{Correlated observables }\Delta g_{\rm NL}\text{ and }\Delta S_8\text{ with slope }A\in[200,500]\text{ (Gate 3)}.
\end{array}
\]

The pathway must hold identically across all three tracks (A, B, C). It is mediated solely by the third cumulant \(C^{(3)}\) (from \(\lambda_6(k)\)) and the running-vacuum coefficient \(\nu\) (from integrated Wetterich flow). The correlation remains falsifiable by its specific exponential functional form.

\section{Expanded Causal Chain}
The full pipeline is shown schematically below (each arrow references the responsible gate and explicit derivation):

\[
\begin{array}{l}
\text{Microscopic UV (AL-GFT / Track B/C)} \\
\quad \downarrow \text{(Gate 1)} \\
\text{Noise kernel } N_{\mu\nu} + C^{(3)} \propto \lambda_6(k_{\rm CMB}) \\
\quad \downarrow \text{(Gates 2 \& 4)} \\
\text{RG flow with controlled truncation} \\
\quad \downarrow \\
\text{Einstein–Langevin backbone} \\
\quad \downarrow \text{(Gate 3)} \\
(\Delta g_{\rm NL},\ \Delta S_8)\ \text{with exponential correlation slope } A \in [200,500]
\end{array}
\]

Explicit steps (with cross-references):

\noindent
\textbf{1. Microscopic foundation (Gate 1, §3):} Arithmetic-Langevin GFT (or Track-B stringy/worldsheet or Track-C inverted multiplicity reconstruction) generates the environmental influence functional. Linear coupling to multiplicity modes produces the leading noise kernel \(N_{\mu\nu}\) and third cumulant \(C^{(3)} \propto \lambda_6(k_{\rm CMB})\).

\noindent
\textbf{2. RG flow and prior derivation (Gate 2, §5–6):} The sextic truncation of \(\Gamma_k[\phi]\) (melonic + first non-melonic) is evolved via the Wetterich equation from \(k = M_P\) to \(k = H_0\), yielding \(\nu_{\rm eff}(M)\) with the log-link coefficient \(c\) distributed as \(\mathcal{N}(1937, 544^2)\) (Gate 2 derived distribution).

\noindent
\textbf{3. Truncation control (Gate 4, §4):} Power counting, Ward identities, EPRL asymptotics, and \(\Delta c/c < 0.04\) error bound ensure higher operators contribute negligibly to \(n_{\rm NG}\) and \(C^{(3)}\).

\noindent
\textbf{4. Einstein–Langevin backbone (Gate 1, §5):} The noise kernel and running parameters enter
\[
G_{\mu\nu}[g] + \Lambda_{\rm eff}(H) g_{\mu\nu} = 8\pi G_{\rm eff} \bigl(\langle T_{\mu\nu}\rangle_{\rm ren} + \xi_{\mu\nu}\bigr),
\]
with \(\Lambda(H) = \Lambda_0 + \nu H^2\) and stochastic tensor \(\xi_{\mu\nu}\) having cumulants fixed by the influence functional.

\noindent
\textbf{5. Observable mapping (Gate 3, §2):} CMB channel: \(C^{(3)}(k_{\rm pivot})\) sources tree-level \(g_{\rm NL}^{\rm eff} \sim C^{(3)}/P_\zeta^2\) (relative to baseline \(g_{\rm NL0}\) in \(\Lambda\)CDM); LSS channel: \(\nu\) shifts the growth factor, producing \(\Delta S_8/S_8 \propto \nu \cdot f(k_{\rm LSS}, z_{\rm eff})\). Eliminating the common \(\lambda_6(k)\) running yields the predicted exponential correlation.

\section{Convergence of the Three Tracks}
All tracks feed the identical Einstein–Langevin backbone (identical noise kernel statistics and \(\nu_{\rm eff}\) injection):
\begin{itemize}
    \item \textbf{Track A} (pure AL-GFT): Zeta-Comb environment.
    \item \textbf{Track B} (string-enhanced): Worldsheet \(\beta_{\rm ws}\) perturbation (bounded by \(\alpha_{\rm ws}\)).
    \item \textbf{Track C} (inverted digital twin): Multiplicity cumulant reconstruction via regularized inverse problem.
\end{itemize}
Track-specific corrections remain sub-percent (as bounded in Gate 4) and shift the universal slope \(A\) by less than \(\sim 7\%\) relative to the central band.

\section{Uniqueness and Falsifiability}
The exponential form---locked by logarithmic RG integration of \(\nu\) versus linear vertex dependence of \(C^{(3)}\)---is structurally distinct from linear/power-law shifts in modified-gravity, neutrino-mass, or dark-energy extensions of \(\Lambda\)CDM. Only the single UV coupling \(\lambda_6(k)\) mediates the link (i.e., no additional parameters can move a point along or off the predicted line). Detection of a steep positive correlation with slope inside \([200, 500]\) (or null case \(\Delta S_8 \approx 0\), \(g_{\rm NL} \lesssim 10^4\)) tests the framework directly. Track C serves as an internal consistency check: reconstruction succeeds only if recovered parameters lie inside the predicted band.

\section{Observational Confrontation Roadmap}
The framework is now ready for:
\begin{itemize}
    \item Joint MCMC sampling over the unified parameter space \(\{M, c, n_{\rm NG}\) (bounded), track-specific nuisance parameters\}.
    \item Confrontation with forthcoming datasets (DESI Y5, LiteBIRD, Euclid) in the \((\Delta S_8, \log g_{\rm NL})\) plane.
    \item Bayesian evidence comparison against \(\Lambda\)CDM and standard extensions.
\end{itemize}

\paragraph{Current status:}
All five gates are validated. The framework now stands as a complete, modularly falsifiable pipeline from microscopic quantum geometry (AL-GFT / string-enhanced / multiplicity-twin reconstructions) to the correlated CMB+LSS observables \((\Delta g_{\rm NL}, \Delta S_8)\). The predicted exponential correlation with slope \(A \in [200,500]\) is ready for confrontation with DESI Y5, LiteBIRD, and Euclid.

\nocite{*}
\bibliographystyle{unsrt}
\bibliography{references}
\clearpage
\end{document}